% !TEX program = pdflatex
%
% Автореферат магистерской диссертации.
% Тем же классом pmk-graduate, что и сама работа, но в режиме [master,summary]:
% класс грузит A5 twoside и сам рисует титульный лист автореферата
% (\maketitlepage). Цель и задачи берутся из общего metadata.tex — ТЕ ЖЕ,
% что и в дипломе (\goalname и список задач \tasksname).
%
% Сборка содержания:        pdflatex autoreferat
% Раскладка «книжкой» (A4): см. autoreferat_booklet.tex
%
% Длинных тире в тексте нет (стилевое требование); em-dash остаётся только
% в служебных строках самого класса и в библиографической ссылке по ГОСТ.
%
\documentclass[master,summary]{pmk-graduate}

\usepackage{amsmath}
\usepackage{amssymb}

\begin{document}
% ООП ФИИТ: Информационные технологии в управлении и принятии решений
% (направление 02.04.02 «Фундаментальная информатика и информационные технологии»)
\program{mitcd}

\title{%
Определение структурных сегментов и эмоциональных характеристик
музыкальных композиций методами машинного обучения%
\english
Structural Segmentation and Emotional Characterisation of Musical
Compositions by Machine Learning Methods%
}

\student{%
Антонов Михаил Евгеньевич%
\brief М.\,Е.\,Антонов%
\group М-26%
\english Antonov Mikhail Evgenyevich%
}

\superviser{%
Кудряшов Максим Юрьевич%
\brief М.\,Ю.\,Кудряшов%
\position доцент кафедры информационных технологий%
\degree кандидат физико-математических наук%
}

\year{2025-2026}

\goal{%
Разработка и экспериментальное исследование программного модуля
автоматического анализа музыкальных композиций, выполняющего
структурную сегментацию, оценку эмоциональных характеристик и
жанровую классификацию аудиозаписей на основе мультизадачной
модели глубокого обучения.%
}

% ---------------------------------------------------------------------------
% Задачи: перечень переписывается в подраздел «Цель и задачи» и повторно
% используется в подразделе «Основные результаты». Формат элемента:
%   \task{<формулировка задачи>\result{<краткий результат>}}
% ---------------------------------------------------------------------------
\begin{tasks}
\task{%
Провести обзор области Music Information Retrieval и обосновать
мультизадачный подход к объединению четырёх задач.%
\result{%
Выполнен обзор 40 источников по трём задачам MIR
и пяти семействам аудио-архитектур (Глава~1). По итогам
обзора в качестве методологической основы выбрана
мультизадачная модель с общим аудио-энкодером и четырьмя
независимыми классификационными головами. Аргументы в её
пользу (содержательная связь задач, регуляризующий эффект
общего ствола и экономия вычислений при инференсе) собраны
в~\textsection\,2.3.%
}}

\task{%
Подготовить обучающий корпус на основе трёх открытых датасетов
(DEAM, Harmonix~Set, GTZAN) с унифицированным пайплайном
извлечения признаков.%
\result{%
Собран объединённый корпус из 3\,545 композиций на основе
трёх открытых датасетов: DEAM (эмоции), Harmonix~Set
(структура, с автоматическим скачиванием 743~треков
с YouTube через \texttt{yt-dlp}) и GTZAN (жанр). Реализован
модуль \texttt{musicml.features} для унифицированного
извлечения логарифмических мел-спектрограмм со 128~фильтрами,
нарезки на десятисекундные окна с шагом 1~секунда и
глобальной нормализации. Итоговый обучающий набор насчитывает
176\,127 окон, разбиение на обучающую, валидационную и
тестовую выборки 70/15/15 делается на уровне треков, что
исключает утечку между сплитами.%
}}

\task{%
Разработать и реализовать пять архитектур мультизадачных
моделей: CNN, CNN\,+\,BiLSTM, PANNs\,+\,Linear,
PANNs\,+\,BiLSTM и AST.%
\result{%
Реализованы и обучены пять архитектур, охватывающих ключевые
оси сравнения. В их числе: свёрточная сеть с блоками
канального внимания (около 800~тысяч параметров), её надстройка
двунаправленной LSTM, линейный классификатор и BiLSTM поверх
замороженных эмбеддингов PANNs CNN14, а также частично
размороженный Audio Spectrogram Transformer
с 14{,}4\,млн обучаемых параметров из 87\,млн общих.
У всех конфигураций единый интерфейс выхода: четыре
классификационные головы плюс две регрессионные для непрерывных
значений возбуждения и валентности. Обучение всех пяти
проведено на одном тестовом корпусе с единым протоколом,
чтобы результаты были напрямую сопоставимы.%
}}

\task{%
Выполнить экспериментальное сравнение пяти архитектур по четырём
задачам на едином тестовом корпусе.%
\result{%
Проведены четыре последовательных эксперимента: базовый
запуск v1, улучшенный v2, расширение трансформер-архитектурой
AST~v2 и неудачный CNN~v3 с агрессивным перевзвешиванием
классов (документированный negative result). Лучший по
среднему результату оказался AST~v2: точность 58{,}5\,\%
усреднённо по четырём задачам, 40{,}6\,\% на структурной
сегментации (на 5{,}3 процентных пункта выше CNN~v2) и
81{,}5\,\% на жанровой классификации. Победители по отдельным
задачам распределились между AST~v2 (сегментация),
PANNs+BiLSTM (возбуждение и валентность) и CNN~v2 (жанр).
Одной архитектуры, доминирующей сразу по всем четырём задачам,
не нашлось.%
}}

\task{%
Реализовать постобработку сегментных предсказаний на базе
декодера Витерби с музыкологическими априорными распределениями.%
\result{%
Разработан и реализован декодер~v2 для постобработки сегментных
предсказаний. В нём объединены шесть механизмов: температурное
масштабирование логитов, асимметричная матрица переходов
6$\times$6 с музыкологическими априори, позиционные приоры
для вступления и концовки, поклассовые ограничения минимальной
длительности, привязка границ к локальным максимумам кривой
новизны и проверка повторяющихся сегментов на близость
представлений. Полный пайплайн повышает Boundary~F1 с 0{,}31
до 0{,}45 (на 14 процентных пунктов выше базового аргмакса)
и сокращает среднее число выдаваемых сегментов
с 28 до 7--8 на трёхминутный трек.%
}}

\task{%
Разработать веб-прототип для интерактивной демонстрации
результатов.%
\result{%
Реализован полнофункциональный веб-прототип с трёхзвенной
архитектурой: фронтенд на React\,19 с Vite, бэкенд на Elysia
поверх Bun, сервис инференса на FastAPI с предобученной AST~v2
в памяти. Центральное решение интерфейса~--- синхронизированный
дашборд из шести панелей с общим временным курсором, в который
сведены разметка структуры, эмоциональные кривые, жанровое
распределение, мел-спектрограмма и сводная панель характеристик.
Дополнительно реализована автоматическая подгрузка
метаданных треков из ID3-тегов и YouTube.%
}}

\task{%
Сформулировать ограничения подхода и направления дальнейшей
работы.%
\result{%
Систематически проанализированы пять групп ограничений текущего
подхода. К ним относятся смещение жанровой головы на современной
музыке из-за устаревшей таксономии GTZAN, фундаментальный потолок
шестиклассовой сегментации в районе 50\,\% по причинам данных
и разметки, размер AST в 574\,МБ с восьмисекундным временем
инференса на трек, невозможность мульти-жанровых предсказаний
в текущей постановке и затруднения с очень длинными треками.
Для каждой группы предложены конкретные направления развития:
надстройка sequence-level декодера BiLSTM или CRF поверх
AST-эмбеддингов, дообучение на современных датасетах
FMA-medium или SALAMI, дистилляция знаний из AST в более
компактный CNN, потоковый инференс со скользящим окном.%
}}
\end{tasks}

\maketitlepage

\section{Актуальность}

Объём музыкальных записей в стриминговых сервисах и цифровых архивах растёт
быстрее, чем возможности их ручной разметки. Из-за этого автоматический анализ
композиций, то есть выделение структурных секций, оценка эмоциональной окраски
и определение жанра, выделился в самостоятельное направление области поиска
музыкальной информации (Music Information Retrieval).

Обычно эти задачи решают порознь, обучая под каждую отдельную модель. Однако они
опираются на общие свойства звука и содержательно связаны между собой: припев и
куплет различаются по энергии и настроению, а смена секции часто совпадает со
сменой эмоции. Это оправдывает мультизадачный подход, при котором один
аудио-энкодер обслуживает сразу несколько классификационных голов. Совместное
обучение даёт регуляризующий эффект, экономит вычисления при работе модели и
позволяет анализировать трек согласованно по нескольким измерениям.

\section{Цель работы}

\goalname

\section{Задачи}

Для достижения поставленной цели в работе решались следующие задачи:
\begin{enumerate}
\let\task\tasktask
\tasksname
\end{enumerate}

\section{Постановка задачи и данные}

Система решает четыре задачи классификации, вычисляемые для каждого окна записи.
Структурная сегментация относит окно к одной из шести секций: вступление, куплет,
припев, мост, инструментальный фрагмент, концовка. Возбуждение и валентность по
модели Рассела оцениваются по трём градациям каждая (низкая, средняя, высокая).
Жанровая классификация работает с десятью категориями коллекции GTZAN.

Входом модели служит логарифмическая мел-спектрограмма со 128 частотными
фильтрами, нарезанная на окна длиной 10 секунд с шагом 1 секунда. Обучающий
материал собран из трёх открытых датасетов: DEAM с эмоциональными аннотациями,
Harmonix Set со структурной разметкой (743 трека загружены с видеохостинга
утилитой yt-dlp) и GTZAN с жанровыми метками. Объединённый корпус насчитывает
3545 композиций и 176\,127 окон; разбиение на обучающую, валидационную и тестовую
выборки в пропорции 70/15/15 выполнено на уровне треков, что исключает утечку
данных между ними.

\section{Мультизадачная модель и архитектуры}

Модель состоит из общего аудио-энкодера и набора независимых голов: четырёх
классификационных по числу задач и двух регрессионных для непрерывных значений
возбуждения и валентности. Обучение ведётся в мультизадачном режиме с чередованием
батчей из разных датасетов и суммарной функцией потерь
\[
\mathcal{L} = \sum_{t} \lambda_t\, \mathcal{L}_t ,
\]
где \(\mathcal{L}_t\) обозначает focal loss отдельной головы, а \(\lambda_t\)
задаёт её вес. Дополнительно применяются сглаживание меток, аугментация
SpecAugment, оптимизатор AdamW с расписанием скорости обучения OneCycle и обрезка
нормы градиента.

Чтобы оценить вклад временного моделирования и предварительного обучения
энкодера, реализованы и обучены пять архитектур:
\begin{itemize}
\item свёрточная сеть с блоками канального внимания, обрабатывающая каждое окно
независимо (около 800 тысяч параметров);
\item та же сеть с надстройкой двунаправленной LSTM, учитывающей соседние окна;
\item линейный классификатор поверх замороженных эмбеддингов PANNs CNN14
размерности 2048;
\item двунаправленная LSTM поверх тех же эмбеддингов;
\item частично дообученный аудио-трансформер AST, предобученный на наборе
AudioSet (14,4 млн обучаемых параметров из 87 млн).
\end{itemize}

\section{Программная реализация}

Модель и весь конвейер обработки оформлены в виде пакета musicml на языке Python
с использованием библиотек PyTorch и librosa. Для наглядной демонстрации построен
веб-прототип с трёхзвенной архитектурой: сервис инференса на FastAPI держит в
памяти веса лучшей модели, серверная часть на Elysia поверх среды Bun управляет
коллекцией треков, а клиент на React выводит результат. Центральным элементом
интерфейса служит синхронизированный дашборд из шести панелей с общим временным
курсором, объединяющий структурную разметку, эмоциональные кривые, жанровое
распределение и мел-спектрограмму.

\section{Апробация и результаты}

Проведены четыре последовательных эксперимента: базовый запуск, его улучшенная
версия, расширение трансформером AST и отдельный неудачный запуск свёрточной сети
с чрезмерным перевзвешиванием классов, зафиксированный как отрицательный результат.
Лучшей по среднему качеству оказалась конфигурация на основе AST: средняя точность
58,5\,\% по четырём задачам, 40,6\,\% на структурной сегментации (на 5,3 процентного
пункта выше лучшей свёрточной модели) и 81,5\,\% на классификации жанра. Победители
по отдельным задачам распределились между AST (сегментация), связкой PANNs с
двунаправленной LSTM (возбуждение и валентность) и свёрточной сетью (жанр);
архитектуры, доминирующей сразу по всем задачам, выявить не удалось.

Для сегментных предсказаний разработан декодер на основе алгоритма Витерби. Он
объединяет температурное масштабирование вероятностей, асимметричную матрицу
переходов с музыкологическими априорными ограничениями, позиционные приоры для
вступления и концовки и поклассовые ограничения минимальной длительности секции.
Постобработка повышает граничную F1-меру с 0,31 до 0,45, то есть на 14 процентных
пунктов, и сокращает среднее число выдаваемых секций с 28 до 7--8 на трёхминутный
трек.

\section{Структура работы}

Работа состоит из введения, семи глав, заключения, списка литературы и приложений.
Первая глава содержит обзор области и обоснование выбранного подхода. Вторая
формализует постановку задачи. Третья глава описывает сбор данных и извлечение
признаков, четвёртая представляет архитектуру модели и процедуру обучения. Пятая
глава посвящена экспериментам и сравнению конфигураций, шестая описывает
веб-прототип, седьмая содержит анализ результатов и ограничений подхода. В
заключении собраны основные выводы и направления дальнейшей работы.

\section{Основные результаты}

\begin{enumerate}
\item Обоснован мультизадачный подход к совместному решению задач структурной
сегментации, оценки эмоциональных характеристик и классификации жанра.
\item Собран объединённый корпус из 3545 композиций на основе датасетов DEAM,
Harmonix Set и GTZAN с единым конвейером извлечения признаков.
\item Реализованы и обучены пять архитектур; лучшей оказалась конфигурация на
основе аудио-трансформера AST со средней точностью 58,5\,\% по четырём задачам.
\item Разработан декодер на основе алгоритма Витерби, повысивший граничную
F1-меру сегментации на 14 процентных пунктов.
\item Построен веб-прототип с синхронизированным дашбордом для интерактивного
анализа произвольных треков.
\item Определены ограничения подхода и намечены направления его дальнейшего
развития.
\end{enumerate}

\section{Публикации автора по теме ВКР}

Антонов М.\,Е. Мультизадачный анализ музыкальных композиций с помощью
предобученного аудио-трансформера // Студенческая научная конференция факультета
прикладной математики и кибернетики: сборник трудов.~--- Тверь: Тверской
государственный университет, 2026.

\end{document}
